\subsection{Stability and Explosion Control in Recursive Generators}
\label{sec:explosion-treatment}

Recursive synthetic generators repeatedly feed their outputs back into the
next evaluation. Feedback can amplify perturbations and produce an explosive
trajectory. This subsection first grounds the implemented Jacobian-based
certificate through contraction, lag-state augmentation, robust invariance,
and random function iterations. It then reviews two additional foundations:
moment contraction for nonlinear autoregressions and a stationary
parameterization of multivariate linear autoregressions. Together they ground
four implemented comparison methods: the original random AST baseline, the
bounded-domain contraction generator, and generators adapted from each of
the two additional papers. Published mathematical results and the
benchmark adaptations are identified separately below.

Four properties guide the interpretation. An invariant domain keeps every
allowed trajectory inside a specified set. A uniform second-moment bound
controls expected squared magnitude over repeated realizations. Contraction
controls how differences between initial histories diminish under shared
future innovations. Strict stationarity means that joint probability laws
are unchanged by a time shift. A method can establish more than one of
these properties; each conclusion below is tied to its stated hypotheses.

\subsubsection{Contraction of Discrete-Time Nonlinear Systems}

\citet{tran2019convergence} compare convergence to one bounded trajectory
with incremental stability, which compares different trajectories.
Their Definition~6 controls small state perturbations through the Jacobian
in a uniformly bounded metric with a uniformly bounded inverse.
For the fixed nonsingular metric transformation $\Theta$ relevant to this
implementation, the transformed Jacobian is
\begin{equation}
 \mathcal J_\Theta(t,\mathbf{s})
 =\Theta J_{\mathcal F}(t,\mathbf{s})\Theta^{-1},
 \qquad
 \mathcal J_\Theta^\top\mathcal J_\Theta\preceq(1-\mu)I,
 \quad \mu>0.
\end{equation}
The transformation specifies how distances are measured. Theorem~15
relates contraction to exponential incremental stability; Theorem~16
adds compact positive invariance to obtain convergence to a unique bounded
trajectory. The benchmark uses a fixed induced norm and a convex invariant
box with $\sup_{\mathbf{s}\in D}\|J_{\mathcal F}(\mathbf{s})\|\leq\gamma<1$.
The derivative bound limits amplification of a difference along the entire
segment between two histories. With the same innovation at each future
step, repeated updates then decrease this difference geometrically.

\subsubsection{Lagged Recurrences as an Augmented State}

\citet{laraba2016setinvariance} review invariance for linear delay-difference
equations

\begin{equation}
    \mathbf x_{t+1}=\sum_{d=0}^{L-1}\mathbf C_{d+1}\mathbf x_{t-d}.
\end{equation}

$\mathbf C_{d+1}$ multiplies the vector at offset $d$; the next-target lag
is $\ell=d+1$. The authors
stack the current and delayed values into one state, so the delayed
equation becomes an ordinary first-order discrete-time system with a companion
transition matrix. They also define a delay-invariant set: if every initial lag
value lies in a set, every value subsequently produced by the recurrence stays
in that set. This supplies the appropriate state representation for an AST
whose leaves access several lags. In the nonlinear generator, the same idea is
used with $\mathbf{s}_t=(x_t,\ldots,x_{t-L+1})^\top$ and the set condition is
written directly as $\mathcal F_{\varepsilon}(D)\subseteq D$. The construction
extends to multivariate recurrences by stacking all $N$ variables at all $L$
lags into a state in $\mathbb R^{NL}$.

\subsubsection{Contraction with Stochastic Innovations}

\citet{kawano2024contraction} study systems
$\mathbf{s}_{t+1}=\mathcal F(\mathbf{s}_t,\xi_t)$.
Their incremental moment stability compares two trajectories under the same
stochastic realization and bounds the expected distance by a geometric
decay. Their Jacobian and metric conditions cover general stochastic
processes, with specializations for temporally independent and Markov
parameters.

The benchmark's bounded-domain method imposes a stronger, pathwise
condition: its Jacobian bound holds for every allowed state and every
innovation satisfying $|\varepsilon_t|\leq b_\varepsilon$.
With the same realized innovations, two histories therefore become
geometrically closer. A stationary long-run law additionally depends on
the update rule and the innovation process. The benchmark supplies a fixed
rule and i.i.d. innovations; the random-function result below provides
the corresponding distributional interpretation.

\subsubsection{Stationary Laws of Random Function Iterations}

\citet[Theorem~1]{diaconis1999iterated} represent a Markov chain as
$\mathbf{s}_{t+1}=\mathcal F_{\xi_{t+1}}(\mathbf{s}_t)$ with i.i.d.
random Lipschitz maps. On a complete separable metric space, the theorem
requires a finite expected Lipschitz factor, a finite expected displacement
from some reference state, and contraction on average:
\begin{equation}
 \mathbb E[\log K_\xi]<0.
\end{equation}
Here $K_\xi$ bounds the amplification of a state difference by the sampled
map. The conclusion is a unique stationary distribution and geometric
convergence in the Prokhorov distance, a distance between probability laws.
The benchmark's uniform bound $K_\xi\leq\gamma<1$, together with the
bounded invariant domain, supplies a stronger sufficient construction.
The result explains why persistent innovations can sustain fluctuations
while the distribution forgets the initial state. Forward simulations
approach a stationary law; coupling two starting histories with the same
noise gives a separate comparison of their realized paths.

\subsubsection{Shao and Wu: Moment Contraction for Nonlinear Time Series}
\label{sec:shao-wu-gmc}

\paragraph{Published result.}
\citet{shao2007spectral} develop spectral theory for nonlinear time series.
Section~5, Theorem~5.1 (pp.~11--12 of the linked electronic reprint), provides
the relevant generation result. With i.i.d. innovations and a measurable,
time-homogeneous recurrence $\mathcal R_{\mathcal T}(\mathbf{s};\varepsilon)$,
its second-moment specialization assumes a finite second moment at some
reference history and the global bound
\begin{equation}
\begin{split}
 \|\mathcal R_{\mathcal T}(\mathbf{s};\varepsilon)
       -\mathcal R_{\mathcal T}(\mathbf{s}';\varepsilon)\|_{L^2}
 &\leq \sum_{d=0}^{L-1}a_d^{\mathrm{GMC}}|s_d-s'_d|,\\
 a_d^{\mathrm{GMC}}\geq0,\qquad
 \sum_{d=0}^{L-1}a_d^{\mathrm{GMC}}&<1.
\end{split}
\label{eq:gmc-generator}
\end{equation}
Here $\|Z\|_{L^2}=(\mathbb E[\|Z\|_2^2])^{1/2}$. The theorem yields a
causal strictly stationary solution and geometric moment contraction (GMC):
the expected squared difference between coupled versions with a shared future
innovation sequence decays geometrically. Vector-valued states are admitted;
the absolute differences become norms of lag blocks. Theorem~5.2 preserves
GMC under stable ARMA filtering, and Examples~5.1--5.5 cover exponential
autoregression, conditional heteroscedasticity, and bilinear recursions.

\paragraph{Benchmark adaptation.}
The standard factor mode wraps lag leaves involved in nonlinear products
with $\tanh$, retains standalone linear terms, and scales the AST until
its global sensitivity budget is below one. For example, the drift
$0.3\tanh(x_t)\tanh(x_{t-1})$ has
derivative bounds $0.3$ for each lag, giving a total bound of $0.6$; Gaussian
innovations satisfy the required moment condition. This route targets
stationary nonlinear dynamics on an unbounded state space. Global derivative
bounds can overestimate sensitivity when repeated factors represent the
same historical coordinate. The alternative standard mode applies $\tanh$
after a complete product and checks the uniform amplitude moment bound
derived in Section~\ref{sec:implemented-shao-wu}. That bound is a project
construction with its own stated conclusion. The paper's GMC conclusion
applies when its global sensitivity and moment hypotheses are verified.
A finite
direct uniform initialization can leave a finite-horizon transient; its law
is recorded, and no preliminary path is simulated or discarded. The AST
sampling and certification rules are project adaptations; deriving compatible
explanation labels from the retained equations is a further benchmark step.

\subsubsection{Heaps: Stationary VAR Models through Reparameterization}
\label{sec:heaps-var}

\paragraph{Published result.}
\citet{heaps2023stationarity} constructs priors supported on stationary Gaussian
VAR models. For $\mathbf{x}_t=\sum_{\ell=1}^{L}\mathbf C_\ell
\mathbf{x}_{t-\ell}+\boldsymbol\varepsilon_t$, causal stationarity requires
the roots of \mbox{$\det(I-\sum_{\ell=1}^{L}\mathbf C_\ell z^\ell)=0$} to lie
outside the unit circle, equivalently $\rho(\mathbf B)<1$ for the companion
matrix $\mathbf B$. Sections~2.1--2.2 (pp.~3--5 of the arXiv version) map
unconstrained matrices to partial-autocorrelation matrices:
\begin{equation}
 \mathbf P_\ell^{\mathrm{PAC}}
 =\left(I+\mathbf U_\ell\mathbf U_\ell^\top\right)^{-1/2}\mathbf U_\ell,
 \qquad \ell=1,\ldots,L.
\label{eq:heaps-pac-map}
\end{equation}
The symmetric inverse square root maps a singular value $u\geq0$ to
$u/\sqrt{1+u^2}<1$. The Appendix (pp.~11--12) reconstructs the VAR coefficients
from these matrices and a positive-definite innovation covariance.
\citet{ansley1986reparameterizing} established the earlier unconstrained
parameterization; Heaps introduces symmetric square roots, a compatible
reverse recursion, and interpretable priors preserving useful symmetries.

\paragraph{Benchmark adaptation.}
The generator takes a specified innovation covariance
$\Sigma_{\varepsilon}\succ0$ and samples $\mathbf U_\ell$. It applies the
reverse recursion and stores the resulting $\mathbf C_\ell$ as linear ASTs.
Sampling the joint lag history from its stationary Gaussian law gives strict
and second-order stationarity from the first generated sample. Signed
coefficients and variable--lag support are read from the final equations.
Sections~3.3--3.4 distinguish shrinkage of intermediate parameters from
structural sparsity: isolated zeros generally change under the transformations.
Accordingly, support must be measured after reconstruction, and an imposed
coefficient mask requires renewed certification. The sampling distribution
also controls persistence and mechanism diversity. This method extends the
benchmark's earlier linear examples to jointly parameterized VAR models
with several variables and lags.

\subsubsection{Implications and Four Generator Strategies}
\label{sec:four-generator-strategies}

The AST families evaluate a retained deterministic expression as

\begin{equation}
    x_{t+1}
    =F_{\mathcal T}(x_t,\ldots,x_{t-L+1})+\varepsilon_{t+1}
\end{equation}

The bounded-domain constructor certifies the induced transition of the
augmented state. For
$D=[-r_D,r_D]^L$, interval propagation supplies both an output bound for
$F_{\mathcal T}$ and derivative bounds
$\overline d_d\ge
\sup_{\mathbf{s}\in D}|\partial F_{\mathcal T}/\partial s_d|$. With
$w_d=q_w^d$ and $q_w>1$, the corresponding weighted-infinity-norm bound is

\begin{equation}
    \gamma
    =\max\!\left\{
      \sum_{d=0}^{L-1}\overline d_d q_w^d,
      q_w^{-1}
    \right\}<1.
\end{equation}

The first term controls the nonlinear recurrence row of the Jacobian; the
second controls the deterministic shift of values through the lag state.
For $L=1$, there is no shift row and only the first term is used.
Bounded innovations are handled separately through the robust-invariance
condition

\begin{equation}
    \sup_{\mathbf{s}\in D}|F_{\mathcal T}(\mathbf{s})|
    +b_{\varepsilon}
    \le r_D,
\end{equation}

which ensures $\mathcal F_{\varepsilon}(D)\subseteq D$ for every admissible
innovation. Contraction therefore certifies exponential forgetting of the
initial history, invariance prevents escape from the certified domain, and the
random-iteration result supports a unique long-run distribution. Monte Carlo
simulation complements both certificates by diagnosing implementation errors
and measuring their conservatism.

These conditions are sufficient and intentionally conservative. A bounded
oscillatory mechanism may fail the uniform Jacobian test, interval arithmetic
may overestimate gains, and a hard invariant domain requires bounded noise.
The certificate consequently defines a controlled subset of the original AST
vocabulary whose trajectories can be generated with an explicit mathematical
stability guarantee.

Table~\ref{tab:four-generator-strategies} summarizes the implemented
comparison. The Shao--Wu factor mode and the bounded-domain method use
contraction with different domains and innovation assumptions. The
product-output mode checks its separately derived uniform second-moment
bound. Heaps provides a linear multivariate
parameterization over the causal stationary region. All four families are
evaluated in Section~\ref{sec:implemented-generator-results}.

\begingroup
\small
\renewcommand{\arraystretch}{1.12}
\begin{longtable}{@{}
>{\raggedright\arraybackslash}p{0.20\textwidth}
>{\raggedright\arraybackslash}p{0.34\textwidth}
>{\raggedright\arraybackslash}p{0.39\textwidth}@{}}
\caption{Four implemented generator strategies. Guarantees apply
under the stated domain, innovation, and initialization assumptions.}
\label{tab:four-generator-strategies}\\
\toprule
Method and status & Mechanism construction & Guarantee and principal trade-off\\
\midrule
\endfirsthead
\toprule
Method and status & Mechanism construction & Guarantee and principal trade-off\\
\midrule
\endhead
\bottomrule
\endlastfoot
Original random AST\newline Implemented &
Sample tree topology, operators, parameters, and lag bindings; execute the
recurrence. &
Empirical baseline with broad sampled algebraic structures; stability and
stationarity are assessed for each mechanism.\\
\addlinespace
Jacobian contraction\newline Implemented &
Scale the AST using interval bounds; certify a weighted lag-state Jacobian
and invariance of a bounded domain. &
Common-noise exponential forgetting and bounded trajectories in the certified
box. Current implementation: scalar series, bounded uniform innovations,
and sums/products.\\
\addlinespace
Shao--Wu family\newline Two standard modes &
Saturate factors and check the global sensitivity budget, or saturate the
complete product and check a growth envelope. &
Factor mode: causal stationary solution and GMC under Theorem~5.1.
Product-output mode: the derived uniform second-moment amplitude bound.
Both use directly sampled histories and Gaussian innovations.\\
\addlinespace
Heaps stationary VAR\newline Implemented &
Sample unconstrained matrices; transform to partial autocorrelations and
reconstruct VAR coefficients. &
Stationary Gaussian VAR($L$) with exact stationary initialization. Linear
interactions; final sparsity and persistence depend on the transformation.\\
\end{longtable}
\endgroup

The comparison reports explosion and near-zero rates together with
mechanism-sampling, certification, initialization, and trajectory-generation
times. Shared seeds, dimensions, lengths, and innovation variance support
matched experiments; innovation laws and supported mechanism families are
also reported. Numerical threshold crossings and finite-sample stationarity
diagnostics are empirical outcomes, while the certificates establish the
declared mathematical properties. Gaussian stationary trajectories can still
produce large excursions.

The benchmark will extract compatible explanation references from each
executable mechanism: variable--lag support, signed linear coefficients,
state-dependent term contributions, and derivatives. Multi-horizon
sensitivities propagate through the full recursion. SHAP references additionally
require a declared coalition and resampling rule.

